\documentclass[11pt,a4paper]{article}

\usepackage[utf8]{inputenc}
\usepackage[T1]{fontenc}
\usepackage{geometry}
\usepackage{booktabs}
\usepackage{longtable}
\usepackage{hyperref}
\usepackage{xcolor}
\usepackage{listings}
\usepackage{amsmath}

\geometry{margin=2.5cm}
\hypersetup{
    colorlinks=true,
    linkcolor=blue,
    urlcolor=blue,
    citecolor=blue
}

\lstdefinestyle{code}{
    basicstyle=\ttfamily\small,
    breaklines=true,
    frame=single,
    columns=fullflexible
}
\lstset{style=code}

\title{MOF\_benchwork: Module A Documentation}
\author{Katharina}
\date{\today}

\begin{document}
\maketitle

\begin{abstract}
This document describes the current state of \texttt{MOF\_benchwork} up to the first runnable version of Module A.
Module A targets single-component CO2 adsorption in a rigid MOF framework using GCMC-style LAMMPS input generation.
The current implementation can resolve benchmark inputs, generate LAMMPS input artifacts, execute pressure-point runs, parse LAMMPS logs, and evaluate the resulting adsorption isotherm against curated reference data.
The simulation output is treated as absolute adsorption from molecule counts, and the evaluation layer can additionally calculate excess adsorption when a pore volume is available.
\end{abstract}

\tableofcontents
\newpage

\section{Project Goal}

The goal of \texttt{MOF\_benchwork} is to provide a reproducible benchmark pipeline for gas adsorption simulations in metal-organic frameworks.
The project is structured around benchmark modules:

\begin{itemize}
    \item Module A: single-component adsorption in a rigid framework.
    \item Module B: mixture adsorption and selectivity.
    \item Module C: Widom insertion, Henry coefficients, or force-field comparison.
    \item Module D: flexible-framework simulations or MD/GCMC flexibility studies.
\end{itemize}

The current implementation focuses on Module A, using MOF-5/IRMOF-1 and CO2 as the first benchmark case.

\section{Current Scope of Module A}

Module A currently covers the following workflow:

\begin{enumerate}
    \item Load a benchmark JSON input file.
    \item Normalize missing configuration values with defaults.
    \item Resolve the material alias MOF-5 to IRMOF-1.
    \item Resolve the CIF file and charge scheme.
    \item Resolve force-field files and adsorbate definitions from CRAFTED.
    \item Select the benchmark module automatically.
    \item Build a side-effect-free run plan.
    \item Materialize a run directory with generated LAMMPS data, molecule templates, force-field include files, and GCMC inputs.
    \item Run a short smoke test or a full pressure-grid isotherm.
    \item Parse LAMMPS logs into JSON summaries.
    \item Evaluate the isotherm against CRAFTED and NIST reference data.
    \item Report both absolute and excess adsorption where possible.
\end{enumerate}



\section{Benchmark Configuration}

The main Module A input is stored in:

\begin{lstlisting}
input_json_files/benchmark_tao2022_mof5_co2.json
\end{lstlisting}

The configuration contains the following conceptual sections:

\begin{itemize}
    \item \texttt{material}: framework name, CIF source, charge scheme, and data source.
    \item \texttt{adsorbates}: adsorbate components, currently CO2.
    \item \texttt{conditions}: temperature and pressure grid.
    \item \texttt{benchmark}: task, metrics, and reference-data configuration.
    \item \texttt{simulation}: LAMMPS/GCMC settings.
    \item \texttt{evaluation}: absolute/excess handling and gas-density model.
    \item \texttt{output}: output directory, run ID, overwrite behavior, and saved artifacts.
\end{itemize}

\subsection{Current Paper-Like Setup}

The current paper-like MOF-5/CO2 setup uses:

\begin{longtable}{p{0.36\textwidth}p{0.48\textwidth}}
\toprule
\textbf{Setting} & \textbf{Current value} \\
\midrule
Material & MOF-5, resolved as IRMOF-1 \\
Charge scheme & DDEC \\
Adsorbate & CO2 \\
Temperature & 298.15 K \\
Pressure points & 0.1, 0.5, 1, 2, 5, 10, 20, 35 bar \\
Force field & UFF framework and CRAFTED CO2 definition \\
Framework model & Rigid \\
Unit cells & 1 x 1 x 1 \\
LAMMPS units & real \\
Pair style & \texttt{lj/cut/coul/long} \\
Cutoff & 12.8 Angstrom \\
KSpace & \texttt{pppm}, accuracy \texttt{1e-4} \\
Fugacity model & Peng-Robinson via CoolProp \\
Production cycles & 500000 \\
Initialization cycles & 500000 \\
\bottomrule
\end{longtable}

\section{Pipeline Design}

\subsection{Planning}

The planning stage is side-effect-free.
It creates a JSON-serializable run plan without writing output files or starting LAMMPS.
The main command is:

\begin{lstlisting}[language=bash]
python benchmark.py benchmark_tao2022_mof5_co2.json --dry-run
\end{lstlisting}

The run plan contains:

\begin{itemize}
    \item selected module,
    \item resolved material and CIF path,
    \item resolved charge scheme,
    \item resolved force-field files,
    \item adsorbate definitions,
    \item pressure and temperature points,
    \item reference candidates,
    \item output paths,
    \item evaluation settings.
\end{itemize}

\subsection{Preparation}

The preparation stage creates a side-effect-free description of all files that should be generated.
The command is:

\begin{lstlisting}[language=bash]
python benchmark.py benchmark_tao2022_mof5_co2.json --prepare
\end{lstlisting}

The actual materialization writes:

\begin{itemize}
    \item framework-only LAMMPS data file,
    \item CO2 molecule template,
    \item LAMMPS force-field include file,
    \item one run-zero input,
    \item one GCMC smoke-test input,
    \item one GCMC input per pressure point,
    \item source copies of CIF, force-field, molecule, and reference files.
\end{itemize}

\section{Converters}

\subsection{CIF to LAMMPS Data}

The CIF converter reads the framework geometry with ASE and separately parses CIF charges from \texttt{\_atom\_site\_charge}.
This is necessary because ASE reads the geometry and cell but does not reliably preserve the charge column as required for LAMMPS \texttt{atom\_style full}.

The generated framework data file contains:

\begin{itemize}
    \item atom count,
    \item atom types,
    \item triclinic box parameters,
    \item masses,
    \item \texttt{Atoms \# full} section with molecule ID, atom type, charge, and Cartesian coordinates.
\end{itemize}

For IRMOF-1 the converter generates a restricted triclinic LAMMPS box.
This preserves the periodic cell geometry in the LAMMPS representation.

\subsection{Molecule Definition to LAMMPS Template}

The molecule converter parses CRAFTED/RASPA \texttt{.def} files for adsorbates such as CO2 and N2.
For CO2 it creates a LAMMPS molecule template with:

\begin{itemize}
    \item three atoms,
    \item two rigid bonds,
    \item global atom type IDs,
    \item partial charges from the pseudo-atom definitions.
\end{itemize}

Because the CO2 template contains bonds, LAMMPS input generation reserves additional per-atom bond and special-neighbor slots using options such as:

\begin{lstlisting}
extra/bond/per/atom 2 extra/special/per/atom 2
\end{lstlisting}

\subsection{Force Field to LAMMPS Include}

The force-field converter parses CRAFTED UFF files and writes LAMMPS pair coefficients.
It assigns atom types for both framework atoms and adsorbate atoms, then applies Lorentz-Berthelot mixing for cross interactions.

\section{LAMMPS Execution}

\subsection{Smoke Test}

A short test run can be started with:

\begin{lstlisting}[language=bash]
conda activate MOF_sim
python benchmark.py benchmark_tao2022_mof5_co2.json --run-test
\end{lstlisting}

The smoke test verifies that LAMMPS can read the generated data, molecule template, and pair coefficients, and that GCMC insertion can occur.

\subsection{Full Isotherm}

A full pressure-grid run can be started with:

\begin{lstlisting}[language=bash]
conda activate MOF_sim
python benchmark.py benchmark_tao2022_mof5_co2.json --run-isotherm --new-run --jobs 4
\end{lstlisting}

The \texttt{--new-run} flag creates a timestamped output directory to avoid overwriting older results.
The \texttt{--jobs} flag controls how many pressure points are executed in parallel.

\section{Evaluation}

\subsection{Absolute Adsorption}

The raw LAMMPS output is converted to an absolute loading using the mean number of adsorbate molecules in the simulation cell:

\begin{equation}
    n_{\mathrm{abs}} =
    \frac{1000 \cdot \langle N_{\mathrm{ads}}\rangle}{m_{\mathrm{framework}}}
\end{equation}

where:

\begin{itemize}
    \item $n_{\mathrm{abs}}$ is the absolute adsorption loading in mol kg$^{-1}$,
    \item $\langle N_{\mathrm{ads}}\rangle$ is the mean number of adsorbate molecules per simulation cell,
    \item $m_{\mathrm{framework}}$ is the framework mass in atomic mass units.
\end{itemize}

The numerical values of mol kg$^{-1}$ and mmol g$^{-1}$ are identical.

\subsection{Excess Adsorption}

The evaluation layer can additionally calculate excess adsorption:

\begin{equation}
    n_{\mathrm{ex}} =
    n_{\mathrm{abs}} - \rho_{\mathrm{gas}} V_{\mathrm{pore}}
\end{equation}

where:

\begin{itemize}
    \item $n_{\mathrm{ex}}$ is the excess adsorption loading in mol kg$^{-1}$,
    \item $n_{\mathrm{abs}}$ is the absolute adsorption loading in mol kg$^{-1}$,
    \item $\rho_{\mathrm{gas}}$ is the bulk gas molar density in mol m$^{-3}$,
    \item $V_{\mathrm{pore}}$ is the pore volume in m$^3$ kg$^{-1}$.
\end{itemize}

For IRMOF-1, the current automatic pore volume source is:

\begin{lstlisting}
CRAFTED-2.0.0/RAC_DBSCAN/CRAFTED_MOF_geometric.csv:AV_cm^3/g
\end{lstlisting}

with:

\begin{equation}
    V_{\mathrm{pore}} = 0.845545~\mathrm{cm^3~g^{-1}}
    = 0.000845545~\mathrm{m^3~kg^{-1}}.
\end{equation}

The gas density is calculated through CoolProp using the configured backend.
For production evaluation, the current backend is:

\begin{lstlisting}
HEOS
\end{lstlisting}

\subsection{Reference Matching}

The evaluation distinguishes reference adsorption bases:

\begin{itemize}
    \item \texttt{simulation\_reference}: compared against absolute simulation loading.
    \item \texttt{absolute}: compared against absolute simulation loading.
    \item \texttt{excess}: compared against excess simulation loading.
    \item \texttt{unknown}: reported as unknown and compared against absolute loading by default.
\end{itemize}

This behavior is controlled by:

\begin{lstlisting}
"reference_matching": "by_basis"
\end{lstlisting}

\subsection{Evaluation Outputs}

For each evaluated run, the pipeline writes:

\begin{itemize}
    \item \texttt{evaluated\_isotherm.csv},
    \item \texttt{evaluation\_report.md},
    \item \texttt{isotherm\_comparison.png}.
\end{itemize}

For comparison against all available references, the pipeline writes:

\begin{itemize}
    \item \texttt{evaluations/reference\_candidates.json},
    \item \texttt{evaluations/combined\_reference\_comparison.csv},
    \item \texttt{evaluations/combined\_reference\_comparison.png},
    \item one evaluation subdirectory per active reference.
\end{itemize}

\section{Reference Data}

The current workflow supports two reference sources:

\begin{itemize}
    \item CRAFTED is used as the primary simulation-reference dataset.
    \item NIST ISODB is used as an additional experimental/literature reference source.
\end{itemize}

NIST entries can be excluded from evaluation if they are detected as outliers or if their adsorption basis is not allowed by the benchmark configuration.
Excluded entries are documented in \texttt{reference\_candidates.json} instead of being silently removed.

\section{Visualization}

LAMMPS trajectory dumps are written for pressure-point runs when dump output is enabled in the generated GCMC input.
The dump files can be opened with the OVITO desktop application.
Coloring by particle type is useful to distinguish framework atom types from inserted CO2 atom types.

Typical dump location:

\begin{lstlisting}
outputs/module_A_co2_isotherm/runs/<run_id>/dumps/gcmc_10bar.lammpstrj
\end{lstlisting}

The number of CO2 molecules changes during GCMC because insertion and deletion moves are part of the grand-canonical ensemble.
This is expected behavior and is the basis for calculating adsorption from the mean molecule count.

\section{Validation Status}

The current test suite checks:

\begin{itemize}
    \item benchmark JSON loading,
    \item default configuration normalization,
    \item MOF-5 to IRMOF-1 alias resolution,
    \item module selection,
    \item CRAFTED molecule parsing,
    \item CIF charge parsing,
    \item LAMMPS data generation,
    \item force-field parsing and pair-coefficient generation,
    \item GCMC input generation,
    \item log parsing,
    \item isotherm CSV generation,
    \item NIST reference parsing,
    \item combined reference evaluation,
    \item excess adsorption comparison against excess references.
\end{itemize}

The latest local test run completed with:

\begin{lstlisting}
Ran 44 tests in 41.431s
OK (skipped=7)
\end{lstlisting}

The skipped tests depend on optional local tools such as ASE, LAMMPS, or the local NIST ISODB checkout.

\section{Current Limitations}

Module A is technically runnable, but the scientific validation is not complete yet.
The most important remaining limitations are:

\begin{itemize}
    \item convergence has not yet been quantified,
    \item uncertainty estimates are not yet calculated,
    \item multiple independent random seeds are not yet automated,
    \item block averaging is not yet implemented,
    \item mean absolute error and mean squared error per reference should be made more explicit in a combined report,
    \item the best-matching reference should be highlighted automatically.
\end{itemize}

\section{Next Steps}

The next development steps for Module A are:

\begin{enumerate}
    \item Add convergence analysis by plotting loading as a function of simulation step or block index.
    \item Add multiple-seed execution for every pressure point.
    \item Calculate statistical uncertainty using standard deviation across seeds or block standard error of the mean.
    \item Refactor the combined evaluation report to show MAE/MSE per reference.
    \item Highlight the best-matching reference dataset.
    \item Add a final Module A report that combines input protocol, simulation results, reference comparison, uncertainty, and plots.
\end{enumerate}

\section{Example Commands}

\subsection{Dry Run}

\begin{lstlisting}[language=bash]
conda activate MOF_sim
python benchmark.py benchmark_tao2022_mof5_co2.json --dry-run
\end{lstlisting}

\subsection{Prepare Files}

\begin{lstlisting}[language=bash]
conda activate MOF_sim
python benchmark.py benchmark_tao2022_mof5_co2.json --prepare --new-run
\end{lstlisting}

\subsection{Run Full Isotherm}

\begin{lstlisting}[language=bash]
conda activate MOF_sim
python benchmark.py benchmark_tao2022_mof5_co2.json --run-isotherm --new-run --jobs 4
\end{lstlisting}

\subsection{Evaluate Existing Run}

\begin{lstlisting}[language=bash]
conda activate MOF_sim
python benchmark.py benchmark_tao2022_mof5_co2.json --evaluate-all --run-id <run_id>
\end{lstlisting}

\section{Appendix: Important Assumptions}

\begin{itemize}
    \item The current first benchmark system is CO2 adsorption in MOF-5/IRMOF-1.
    \item The framework is treated as rigid.
    \item The current framework charge scheme is DDEC.
    \item CRAFTED is treated as an external raw-data source and is not committed as curated project data.
    \item The LAMMPS output molecule count is interpreted as absolute adsorption.
    \item Excess adsorption requires a pore volume and a bulk gas density model.
    \item The current pore volume for IRMOF-1 is taken from CRAFTED geometric descriptors.
\end{itemize}

\end{document}
